\section{Verify a small piece against a fake environment \hfill \live{}}

\textbf{The idea in one sentence.} The property \texttt{A[] not (P1.cs \&\& P2.cs)} mentions two
processes. We currently check it on all ten. Replace the other eight by a single small automaton
that is \emph{at least as permissive} as they are, and check the property on that instead.

If the replacement is genuinely at least as permissive, then anything bad that the real system can
do, the fake one can also do. So if the small system is safe, the big one is safe. The reverse is
not guaranteed --- the fake environment may invent bad behaviour that the real one cannot produce
--- which is why this is a one-way test: a ``safe'' answer is trustworthy, an ``unsafe'' answer
means we learned nothing and must fall back.

\begin{theorem}[Environment abstraction]
Let $S = C \parallel E$, where $C$ contains every component named in the safety property $\varphi$
together with anything sharing state with them, and $E$ is the rest. Let $\hat{E}$ be any timed
automaton with $E \sqsubseteq \hat{E}$ in the ready-simulation sense on the interface alphabet.
Then
\[
  C \parallel \hat{E} \models \varphi \;\Longrightarrow\; C \parallel E \models \varphi .
\]
\end{theorem}

The proof is short because someone else did the hard part: $\sqsubseteq$ is a \emph{precongruence}
for \uppaal{}'s composition operator, including shared variables and committed locations
(Berendsen \& Vaandrager, FORMATS 2008). Precongruence means the relation survives composition, so
$C \parallel E \sqsubseteq C \parallel \hat E$ follows immediately, and $\varphi$ only reads $C$,
which both sides have identically.

\paragraph{What it would buy.} On \texttt{fischer} at $N=10$, taking $\hat E$ to be one generic
replica plus the shared variable abstracted to $\{0, 1, 2, \text{other}\}$ gives roughly $256$
discrete states with $3$ clocks, against $260\,998$ with $10$ clocks. That is a factor of about
$1000$, and --- more importantly --- the abstract system does not grow with $N$ at all.

\paragraph{Is this new? No, and that is the interesting part.} The technique is quotienting, from
Andersen (LICS 1995) and Laroussinie \& Larsen (CONCUR 1995), and it was demonstrated on
\emph{exactly this model} thirty years ago: Kristoffersen, Laroussinie, Larsen, Pettersson and Yi,
TAPSOFT 1997, verified Fischer at $N = 50$ in $172$ seconds using the tool CMC, at a time when
\uppaal{} gave up around $N = 6$. Separately, Abdulla et al.\ (Nordic J.\ Computing, 2007) settled
Fischer for \emph{all} $N$ at once in $2.12$ seconds.

\paragraph{So why is nobody doing it?} Four reasons, all from the sources:

\begin{enumerate}
\item \textbf{CMC no longer exists}, and no successor tool was found.
\item \textbf{The decomposition order was chosen by hand.} The 1997 paper says so itself: the
      optimal order in which to factor out components ``needs a better understanding''. Automating
      that choice is the open problem.
\item \textbf{CMC could not handle \uppaal{} integer variables.} It modelled the shared variable as
      an explicit $(n{+}1)$-location automaton. Real models use bounded integers.
\item \textbf{The parameterized route has a sharp decidability cliff.} Reachability is decidable
      for unboundedly many identical \emph{one-clock} processes; with a controller and several
      clocks per process it is undecidable, and no cutoff theorem covers \uppaal{}-style shared
      bounded integers or binary channels.
\end{enumerate}

\paragraph{Honest risk.} The theory is proven and old. The gap is \emph{automation} --- constructing
$\hat E$ without a human. That is exactly the part nobody has done, so this lead is not ``implement
a known result'', it is ``attempt the part that stalled''. The consolation is that failure is cheap:
one abstract run that is independent of $N$, then fall back to the normal search.
